% ============================================
% DOC-05
% Gerado de molambudos.md — build_miolo.py
% ============================================
\section*{\textbf{DOC-05} --- Relatório de Sessão II}
\markboth{DOC-05 --- Relatório de Sessão II}{}

\noindent \textit{Dr. Heitor Oliveira --- 21 de março de 1979}

\noindent \textbf{Dr. Heitor Oliveira --- 21 de março de 1979}

\vspace{0.5em}\noindent\rule{\textwidth}{0.5pt}\vspace{0.5em}

\noindent \textbf{Paciente:} Joaquim Antônio Correia --- Paciente 1.260
\textbf{Tipo:} Sessão de avaliação psiquiátrica
\textbf{Duração:} 90 minutos (estendida)

\vspace{0.5em}\noindent\rule{\textwidth}{0.5pt}\vspace{0.5em}

\noindent \textbf{Observações do Paciente:}

\noindent O paciente apresentou-se diferente hoje. Mais animado. Quase eufórico. Sorria quando entrei na sala --- um sorriso assimétrico, puxado para o lado direito do rosto, onde o olho amarelo brilha mais intensamente do que na sessão anterior.

\noindent Ele perguntou: \textit{\textquotedblleft{}O senhor sonhou?\textquotedblright{}}

\noindent Respondi que sim, que tive um sonho incomum. Ele não pediu detalhes. Apenas acenou com a cabeça, como se já soubesse.

\noindent Nesta sessão, Joaquim falou continuamente por 75 minutos. Praticamente sem interrupção. Descreveu sua infância no Ceará, a seca de 1915, a caminhada de Senador Pompeu em direção a Fortaleza, a morte do pai e o confinamento no Campo do Alagadiço. Os relatos são vívidos --- há detalhes sensoriais (cheiros, temperaturas, texturas) que são incomuns mesmo em pacientes com memória preservada.

\noindent O tom do relato mudou quando descreveu o \textquotedblleft{}encontro\textquotedblright{}.

\noindent Segundo Joaquim, ele encontrou uma criatura em uma vala de corpos atrás do campo de concentração. A descrição continha elementos psicóticos (olhos amarelos, espinhos de osso, costelas externas), mas a \textit{consistência interna} é notável. Ele descreveu o encontro com a mesma convicção com que descreveu a morte do pai.

\noindent \textbf{Registro textual:}

\noindent \textit{\textquotedblleft{}Ela disse: 'Você também tem fome.' Eu disse que sim. Ela disse: 'Eu como por você. Eu como por todos que não podem comer.' E eu sentei. Sentei na beira da vala e olhei ela comer. Não senti medo. Senti fome. Uma fome que nunca mais passou.\textquotedblright{}}

\noindent \textbf{Análise:}

\noindent - A \textquotedblleft{}criatura\textquotedblright{} parece ser uma manifestação do trauma infantil de Joaquim, possivelmente uma personificação da fome e da morte que testemunhou na infância.
- O encontro na vala pode representar um momento dissociativo: Joaquim pode ter testemunhado algo insuportável (a vala de corpos, a morte de outros) e sua mente criou uma narrativa alternativa para processar o evento.
- A \textquotedblleft{}fome que nunca passa\textquotedblright{} é consistente com o quadro de Transtorno de Estresse Pós-Traumático, onde o evento traumático é revivido indefinidamente.

\noindent \textbf{Plano Terapêutico:}

\noindent - Solicitei acesso ao diário do paciente, que foi negado inicialmente, mas o paciente concordou em me entregar \textquotedblleft{}quando estiver pronto\textquotedblright{}.
- Manter sessões semanais.

\noindent \textbf{Nota Pessoal:}

\noindent Comecei a ter sonhos recorrentes. Não são pesadelos comuns --- são extremamente vívidos, com cheiros e sensações táteis. Acordo com o gosto de terra na boca. Não mencionei isso a Joaquim, mas ele sabe. Ele sabe.

\vspace{0.5em}\noindent\rule{\textwidth}{0.5pt}\vspace{0.5em}

\noindent \textit{Dr. Heitor Oliveira}
\textit{Psiquiatra --- Ala de Segurança Máxima}
\textit{Hospital Colônia de Barbacena}

\vspace{0.5em}
\noindent\textit{\textbf{↪ Rotas:} \rota{MEM-15} \rotasep \rota{DOC-07} \rotasep \rota{LUC-03}}